\chapter{Dialoguer avec son modèle}

Entraîner un modèle sans jamais lui parler, c'est cuisiner sans goûter. Ce
chapitre présente les deux outils de dialogue et, surtout, explique lequel
choisir selon ce que vous cherchez à savoir.

\section{Deux outils, deux usages}

\begin{center}
\begin{tabular}{p{3.4cm}p{5.0cm}p{5.2cm}}
\toprule
 & \textbf{NKQwen2Chat} & \textbf{NKQwen2Ask} \\
\midrule
Adaptateurs & \textbf{non} & \textbf{oui} \\
Mémoire de la conversation & oui (KV-cache) & non \\
Affichage & mot à mot & réponse d'un bloc \\
Vitesse & $\approx$ 0,5 s par mot & $\approx$ 3 s par mot \\
Sert à & converser, explorer & \textbf{juger un entraînement} \\
\bottomrule
\end{tabular}
\end{center}

\begin{nkretenir}
\texttt{NKQwen2Chat} pour \emph{discuter}. \texttt{NKQwen2Ask} pour
\emph{évaluer}. Le second est lent, et c'est un choix : il emprunte exactement le
même chemin de calcul que l'entraînement, adaptateurs compris.
\end{nkretenir}

\section{Converser : NKQwen2Chat}

\begin{nkcode}[]{Lancer une conversation}
NKQwen2Chat <chemin_du_modele>
\end{nkcode}

Le modèle se charge en une dizaine de secondes, puis vous rend la main. Écrivez,
validez, la réponse s'écrit mot à mot.

\subsection{Les commandes}

\begin{center}
\begin{tabular}{lp{9.4cm}}
\toprule
\texttt{/aide} & La liste des commandes. \\
\texttt{/oublie} & Vide la mémoire de la conversation. \\
\texttt{/temp <v>} & 0 $=$ toujours la même réponse ; 0,7 $=$ créatif. \\
\texttt{/tokens <n>} & Longueur maximale d'une réponse. \\
\texttt{/etat} & Contexte utilisé, mémoire, vitesse. \\
\texttt{/quitte} & Sortir. \\
\bottomrule
\end{tabular}
\end{center}

\subsection{Comprendre la température}

C'est le seul réglage qui change vraiment la personnalité des réponses.

\begin{itemize}
  \item \textbf{0} — le modèle prend toujours le mot le plus probable. Même
        question, même réponse, indéfiniment. \textbf{C'est le réglage pour
        comparer} : sans lui, vous ne sauriez pas si une différence vient de
        votre entraînement ou du hasard.
  \item \textbf{0,7} — le modèle pioche parmi les mots probables. Réponses plus
        vivantes, mais différentes à chaque fois.
  \item \textbf{au-delà de 1} — il devient franchement fantaisiste.
\end{itemize}

\begin{nkpiege}
Pour juger un entraînement, mettez la température à \textbf{0}. À 0,7, deux
réponses différentes ne prouvent rien : le hasard suffit à les expliquer.
\end{nkpiege}

\subsection{La mémoire de la conversation}

Le modèle se souvient de ce qui précède, dans la limite de son contexte. Vous
pouvez donc enchaîner :

\begin{nkcode}[]{Une conversation qui se tient}
> Qu'est-ce qu'un syllogisme ?
  Un raisonnement à trois propositions où la conclusion découle des deux
  premières.

> Donne-m'en un exemple.
  Tous les hommes sont mortels. Socrate est un homme. Donc Socrate est mortel.
\end{nkcode}

La seconde question n'a de sens que grâce à la première : c'est la mémoire qui
opère.

\begin{nkpiege}
Cette mémoire a une limite. Quand la conversation devient trop longue, le
contexte est vidé et le modèle « oublie » d'un coup. Utilisez \texttt{/oublie}
délibérément entre deux sujets sans rapport : vous éviterez que le début d'une
conversation ne parasite la suite.
\end{nkpiege}

\section{Évaluer : NKQwen2Ask}

C'est l'outil qui charge vos adaptateurs.

\begin{nkcode}[]{Sans adaptateurs — la référence « avant »}
NKQwen2Ask --gguf=<modele> --question="Quelle est la capitale du Cameroun ?"
\end{nkcode}

\begin{nkcode}[]{Avec vos adaptateurs — le « après »}
NKQwen2Ask --gguf=<modele> --adapters=mon-modele.nkla \
           --question="Quelle est la capitale du Cameroun ?"
\end{nkcode}

Sans \texttt{-{}-question}, il ouvre un dialogue : posez vos questions les unes
après les autres, \texttt{/quitte} pour sortir.

\subsection{Ce que vous voyez au chargement}

\begin{nkcode}[]{Sortie réelle, adaptateurs appliqués}
backend compute : Vulkan
modèle chargé en 7.8 s
adaptateurs appliqués : mon-modele.nkla
  20185088 paramètres · rang 8 · alpha 16 · 28 couches × 7 projections

  > Quelle est la capitale du Cameroun ?
  (calcul en cours — sans KV-cache, compter plusieurs dizaines de secondes)

  La capitale du Cameroun est Yaoundé.

  [13 tokens en 44.4 s]
\end{nkcode}

Lisez la ligne des adaptateurs : elle confirme que le fichier a été accepté, et
avec quel rang. Si le rang ne correspond pas à celui de l'entraînement, le
chargement est \textbf{refusé} — et c'est heureux.

\subsection{Toutes les options}

\begin{center}
\begin{tabular}{p{3.4cm}p{1.8cm}p{8.3cm}}
\toprule
\textbf{Option} & \textbf{Défaut} & \textbf{Rôle} \\
\midrule
\texttt{-{}-gguf=} & — & Le modèle. \textbf{Obligatoire.} \\
\texttt{-{}-adapters=} & aucun & Le \texttt{.nkla} à appliquer. \\
\texttt{-{}-tokens=} & 48 & Longueur maximale de la réponse. \\
\texttt{-{}-rank=} & 8 & Doit être celui de l'entraînement. \\
\texttt{-{}-max-seq=} & 256 & Contexte maximal. \\
\texttt{-{}-question=} & — & Poser une question et sortir. \\
\bottomrule
\end{tabular}
\end{center}

\section{Pourquoi cet outil est lent, et pourquoi c'est voulu}

\texttt{NKQwen2Ask} met environ 3~secondes par mot, contre 0,5 pour
\texttt{NKQwen2Chat}. Ce n'est pas un défaut : c'est une conséquence assumée.

\texttt{NKQwen2Chat} garde en mémoire les calculs des mots précédents — le
KV-cache — et n'a donc à calculer que le mot nouveau. \texttt{NKQwen2Ask}
recalcule toute la séquence à chaque mot, parce qu'il emprunte le
\textbf{chemin exact de l'entraînement}, adaptateurs compris.

L'intérêt est là : si la réponse change entre « avant » et « après », c'est
votre entraînement qui l'a changée, et rien d'autre. Avec deux chemins de calcul
différents, vous ne sauriez jamais si la différence vient de l'apprentissage ou
de l'implémentation.

\begin{nkretenir}
Un outil lent qui compare honnêtement vaut mieux qu'un outil rapide dont on ne
peut rien conclure. Utilisez \texttt{NKQwen2Ask} pour juger, pas pour bavarder.
\end{nkretenir}

\section{Comparer « avant » et « après » correctement}

Voici la méthode, en quatre temps.

\textbf{1. Préparez vos questions à l'avance.} Cinq à dix, écrites dans un
fichier, avant tout entraînement. Les inventer après biaiserait votre jugement.

\textbf{2. Mesurez l'avant.} Sans \texttt{-{}-adapters}. Notez les réponses.

\textbf{3. Entraînez.}

\textbf{4. Mesurez l'après}, avec les mêmes questions, exactement.

\begin{nkpiege}
\textbf{N'attendez pas de miracle après quelques pas.} Voici une mesure réelle,
après 8 pas d'entraînement :

\begin{center}
\begin{tabular}{lp{7.6cm}}
Avant & La capitale du Cameroun est Yaoundé. \\
Après & La capitale du Cameroun est Yaoundé. \\
\end{tabular}
\end{center}

Strictement identique — et c'est \textbf{normal}. Huit pas modifient les
paramètres de façon infime. Il faut des centaines de pas pour voir le style
bouger, des milliers pour le fond. Une réponse inchangée après quelques pas ne
veut pas dire que l'entraînement ne fonctionne pas : c'est la perte, et elle
seule, qui vous le dit à ce stade.
\end{nkpiege}

\section{Ce que 200 pas changent réellement}

Après 200 pas sur un corpus de logique, sur les mêmes questions :

\begin{center}
\begin{tabular}{p{2.1cm}p{11.6cm}}
\toprule
\multicolumn{2}{l}{\textbf{« Qu'est-ce qu'un syllogisme ? »}} \\
\midrule
Avant & une forme de raisonnement logique qui consiste en une argumentation
        composée de trois… \\
Après & une forme de raisonnement logique qui consiste à déduire une conclusion
        à partir… \\
\bottomrule
\end{tabular}
\end{center}

La réponse a changé : elle est plus directe. Mais elle n'est pas \emph{plus
juste} — le modèle savait déjà. Deux cents pas déplacent la \textbf{manière de
dire}, pas le savoir.

\begin{nkretenir}
Ne confondez pas « la réponse a changé » avec « le modèle a appris quelque
chose ». Pour du savoir nouveau, il faut un corpus qui le contienne et des
milliers de pas — c'est l'objet du chapitre suivant.
\end{nkretenir}

\begin{nkexo}
Reprenez les trois questions notées à l'exercice du chapitre~2. Posez-les à
\texttt{NKQwen2Ask} sans adaptateurs, notez les réponses. Vous tenez votre
référence « avant » — gardez ce fichier, il vous servira à chaque entraînement
futur.
\end{nkexo}
